\documentclass[11pt]{article}
\usepackage[a4paper,margin=18mm]{geometry}
\usepackage{fontspec}
\setmainfont{Latin Modern Roman}
\usepackage{amsmath,amssymb,amsthm,amscd,graphicx,color}
\usepackage{xurl}
\usepackage[unicode,hidelinks]{hyperref}
\allowdisplaybreaks
\setlength\emergencystretch{3em}

\makeatletter
\newtheorem{theorem}{Theorem}
\newtheorem*{theorem*}{Theorem}
\newtheorem{lemma}{Lemma}
\newtheorem*{lemma*}{Lemma}
\newtheorem{corollary}{Corollary}
\newtheorem*{corollary*}{Corollary}
\newtheorem{proposition}{Proposition}
\newtheorem*{proposition*}{Proposition}
\newtheorem{conjecture}{Conjecture}
\newtheorem*{conjecture*}{Conjecture}
{\theoremstyle{definition} \newtheorem{definition}{Definition}
\newtheorem*{definition*}{Definition}
\newtheorem{example}{Example}
\newtheorem*{example*}{Example}
\newtheorem{remark}{Remark}
\newtheorem*{remark*}{Remark}
\newtheorem{note}{Note}
\newtheorem*{note*}{Note}
}
\makeatother
\numberwithin{equation}{section}
\numberwithin{theorem}{section}
\numberwithin{proposition}{section}
\numberwithin{lemma}{section}
\numberwithin{corollary}{section}
\numberwithin{definition}{section}
\numberwithin{example}{section}
\numberwithin{remark}{section}
\numberwithin{note}{section}


\newcommand{\R}{\mathbb R}
\newcommand{\N}{\mathbb N}
\newcommand{\C}{\mathbb C}
\newcommand{\Z}{\mathbb Z}
\newcommand{\T}{\mathbb T}
\newcommand{\eps}{\varepsilon}
\newcommand{\dirint}[2]{\sideset{}{^\oplus}\int\limits_{#1\ }^{\ \ #2}}
\renewcommand{\labelenumi}{(\roman{enumi})}
\newcommand{\rphis}[5]{\,_{#1}\varphi_{#2} \left( \genfrac{.}{.}{0pt}{}{#3}{#4}
\ ;#5 \right)}
\newcommand{\su}{\mathfrak{su}}
\newcommand{\U}{\mathcal U}
\newcommand{\Res}[1]{\underset{#1}{\mathrm{Res}}}
\newcommand{\vect}[2]{\left( \genfrac{.}{.}{0pt}{}{#1}{#2} \right) }


\begin{document}
\begin{center}\Large Unitary decomposition of the specified discrete-pair plus strange-pair coproduct representation of U q su(1,1) when k1 and k2 exceed one half and their difference has absolute value at most one half and zero less than q less than one\end{center}
\noindent\textbf{Stable ID:} 2645\par
\noindent\textbf{Authors:} Wolter Groenevelt\par
\noindent\textbf{Original work:} Quantum Analogs of Tensor Product Representations of su(1,1)\par
\noindent\textbf{Version:} Official SIGMA 7 (2011) article 077, DOI 10.3842/SIGMA.2011.077; journal PDF and native TeX acquired 2026-10-01, exact bytes pinned by SHA256\par
\noindent\textbf{Selected task scope:} Theorem3.1(i):0<q<1,k1,k2>1/2,absolute value(k1-k2)<=\allowbreak{}1/2,epsilon real. T is exactly(negative-discrete pi\_k1 tensor positive-discrete pi\_k2) direct sum(strange pi\_(k1-1/2,-epsilon-k1) tensor strange pi\_(k2-1/2,epsilon+\allowbreak{}k2)),composed with the source coproduct. T is unitary equivalent to the principal-series direct integral rho in[0,-pi/(2lnq)] with epsilon prime=\allowbreak{}k2-k1 plus strange-series summands sigma\_j=\allowbreak{}k1+\allowbreak{}k2+\allowbreak{}epsilon+\allowbreak{}1/2+\allowbreak{}j for integer j with sigma\_j>0, using the source normalized bigqJacobi-function intertwiner.\par
\noindent\textbf{Difficulty:} H2: The representation combines a half-line discrete-series sector with a bilateral strange-series sector, and its Casimir has a nontrivial continuous-plus-discrete spectral decomposition. The author constructs the big q-Jacobi-function kernel separately on each weight space, matches the full tridiagonal Casimir coefficients, and assembles the normalized continuous and mass-point transforms. Joint Casimir and weight spectra then identify the precise irreducible principal and strange summands, rather than treating the representation as a single ordinary tensor product.\par
\noindent\textbf{Editorial boundary note (not author text):} The full original balanced-regime Theorem3.1(i), including its continuous-plus-strange summands and exact normalized intertwiner, is retained for k1,k2 greater than one half and absolute difference at most one half. The two coproduct tensor sectors, every weight-space parameter branch, big q-Jacobi measure/basis/recurrence, complete Casimir match and joint weight/Casimir identification are bundled. Published orthonormal-basis and irreducible-characterization results remain cited prerequisites; the author analogous negative-weight calculation and phase-factor/contiguous-relation note are preserved. Complementary/discrete boundary cases and other tensor-sector theorems are not selected.\par
\noindent\textbf{Source:} \url{https://sigma-journal.com/2011/077/}\quad DOI: \url{https://doi.org/10.3842/SIGMA.2011.077}\par
\noindent\textbf{License and attribution:} Copyright retained by the named authors. Published by SIGMA under CC BY 4.0: \url{https://creativecommons.org/licenses/by/4.0/}. Source: \url{https://sigma-journal.com/about.html}. Adaptation consists of standalone excerpt formatting, cover and numbering restoration; original author language and mathematical text are retained.\par
\noindent\textbf{Editorial symbol notice (not author text):} This is a direct sum of two coproduct tensor sectors, not a single tensor product of irreducible quantum representations. Standard bases, signs, weight-space parameter branches and continuous/discrete normalization are all retained. The negative-weight coefficient computation is the source analogous routine case, not an added proof.\par
\medskip\hrule\medskip\noindent\textbf{Author text begins below.}\par
\setcounter{section}{1}
\setcounter{subsection}{0}
\setcounter{subsubsection}{0}
\setcounter{equation}{0}
% Original source lines 105--377
\section[The quantum algebra $\U_q(\su(1,1))$]{The quantum algebra $\boldsymbol{\U_q(\su(1,1))}$} \label{sec:Uq}

The quantized universal enveloping algebra $\mathcal U_q = \U_q\big(\su(1,1)\big)$ is the unital, associative, complex algebra generated by $K$, $K^{-1}$, $E$, and $F$, subject to the relations
\begin{gather}
K K^{-1} = 1 = K^{-1}K, \qquad KE = qEK, \qquad KF= q^{-1}FK, \nonumber\\
 EF-FE =\frac{K^2-K^{-2}}{q-q^{-1}}.\label{eq:commutationrel}
\end{gather}
The Casimir element
\begin{gather} \label{eq:Casimir}
\Omega = \frac{q^{-1} K^2 +qK^{-2}-2}{(q^{-1}-q)^2} + EF= \frac{q^{-1}K^{-2}+qK^2-2}{(q^{-1}-q)^2} +FE
\end{gather}
is a central element of $\U_q$. The algebra
$\U_q$ is a Hopf $*$-algebra with comultiplication $\Delta$ def\/ined on the generators by
\begin{alignat}{3}
& \Delta(K) = K \otimes K,\qquad && \Delta(E)= K \otimes E + E \otimes K^{-1},& \nonumber\\
& \Delta(K^{-1}) = K^{-1} \otimes K^{-1},\qquad && \Delta(F) = K \otimes F + F \otimes K^{-1} ,& \label{eq:comult}
\end{alignat}
and $\Delta$ is extended to $\mathcal U_q$ as an algebra homomorphism. In particular, it follows from~\eqref{eq:Casimir} and~\eqref{eq:comult} that
\begin{gather}
\Delta(\Omega) = \frac{1}{(q^{-1}-q)^{2}} \big[ q^{-1}\big( K^2 \otimes K^2 \big)+ q
\big(K^{-2} \otimes K^{-2}\big) -2 (1\otimes 1) \big] \nonumber\\
\phantom{\Delta(\Omega) =}{} + K^2 \otimes EF +KF\otimes
EK^{-1} + EK \otimes K^{-1}F + EF \otimes K^{-2}.\label{eq:De(OM)}
\end{gather}
The $*$-structure on $\U_q$ is def\/ined on the generators by
\[
K^*=K, \qquad E^*=-F, \qquad F^* = -E, \qquad (K^{-1})^* = K^{-1}.
\]
Note that the Casimir element is self-adjoint in $\U_q$, i.e.~$\Omega^*=\Omega$.

There are, besides the trivial representation, f\/ive classes of irreducible $*$-representations of~$\U_q$, see \cite[Proposition~4]{MMNNSU}, \cite[Section~6]{BK93}. The representations are given in terms of unbounded operators on $\ell^2(\N)$ or $\ell^2(\Z)$. As common domain we take the f\/inite linear combinations of the standard orthonormal basis vectors $e_n$. The representations are unbounded $*$-representations in the sense of Schm\"udgen \cite[Def\/inition~8.1.9]{Sch}. Below we list the actions of the generators $K$, $K^{-1}$, $E$, $F$ on basis vectors $e_n$. The Casimir element~$\Omega$ plays an important role in this paper, therefore we also list the action of $\Omega$.

The representation listed below are completely characterized, up to unitary equivalence, by the actions of~$K$ and~$\Omega$. Let us brief\/ly describe how these actions determine the actions of~$E$ and~$F$ up to a phase factor. Let $\pi$ be a $\U_q$-representation acting on $\ell^2(\Z)$ with orthonormal basis $\{e_n\}_{n \in \Z}$, and suppose that $\pi(K)e_n=q^{n+\eps}e_n$ and $\pi(\Omega)e_n = \omega e_n$ for all $n \in \Z$, where $\eps$ and $\omega$ are real numbers. The commutation relations~\eqref{eq:commutationrel} imply that $\pi(E)e_n = c_n e_{n+1}$ and $\pi(F) e_n = d_n e_{n-1}$ for certain numbers~$c_n$ and~$d_n$. Furthermore, the relation $E^*=-F$ implies $c_n=-\overline{d_{n+1}}$, so that $\pi(FE) e_n = - |c_n|^2 e_n$. On the other hand, from~\eqref{eq:Casimir} and the actions of $K$ and $\Omega$ it follows that
\[
\pi(FE) e_n = \left(\omega-\frac{q^{2n+2\eps+1}+q^{-2n-2\eps-1}-2}{(q^{-1}-q)^2}\right)e_n.
\]
So $c_n$ can be determined up to a phase factor.

The f\/ive classes of irreducible $*$-representations of $\U_q$ are the following:

\textbf{Positive discrete series.} The positive discrete series $\pi^+_k$
are labeled by $k>0$. The representation space is $\ell^2(\N)$
with orthonormal basis $\{e_n\}_{n \in \N }$. The action is
given by
\begin{gather}
\pi^+_k(K)  e_n  = q^{k+n}  e_n, \qquad \pi^+_k\big(K^{-1}\big)  e_n =q^{-(k+n)} e_n, \nonumber \\
\big(q^{-1}-q\big)\pi^+_k(E)  e_n  = q^{-\frac12-k-n} \sqrt{(1-q^{2n+2})(1-q^{4k+2n})}  e_{n+1},  \nonumber\\
\big(q^{-1}-q\big)\pi^+_k(F)  e_n  = -q^{\frac12-k-n}
\sqrt{(1-q^{2n})(1-q^{4k+2n-2})}  e_{n-1},\nonumber\\
\big(q^{-1}-q\big)^2 \pi^+_k(\Omega)  e_n  = \big(q^{2k-1}+q^{1-2k}-2\big)  e_n.\label{eq:pos}
\end{gather}

\textbf{Negative discrete series.} The negative discrete series
$\pi^-_k$ are labeled by $k>0$. The representation space is
$\ell^2(\N)$ with orthonormal basis $\{e_n\}_{n \in \N}$. The action is given by
\begin{gather}
\pi^-_k(K) e_n = q^{-(k+n)} e_n, \qquad \pi^-_k\big(K^{-1}\big) e_n =q^{k+n} e_n,
 \nonumber\\
 \big(q^{-1}-q\big)\pi^-_k(E)  e_n  = -q^{\frac12-k-n}
\sqrt{(1-q^{2n})(1-q^{4k+2n-2})}  e_{n-1},
\nonumber\\
\big(q^{-1}-q\big)\pi^-_k(F) e_n = q^{-\frac12-k-n}
\sqrt{(1-q^{2n+2})(1-q^{4k+2n})} e_{n+1}, \nonumber\\
\big(q^{-1}-q\big)^2 \pi^-_k(\Omega) e_n = \big(q^{2k-1}+q^{1-2k}-2\big) e_n.\label{eq:neg}
\end{gather}

\textbf{Principal unitary series.} The principal unitary series
representations $\pi^{\mathrm P}_{\rho,\eps}$ are labeled by $0\leq\rho <
-\frac{\pi}{2\ln q}$ and $\eps \in [0,1)$, where $(\rho,\eps)
\neq (0,\frac12)$. The representation space is $\ell^2(\Z)$ with
orthonormal basis $\{e_n\}_{n \in \Z}$. The action is given by
\begin{gather}
\pi^{\mathrm P}_{\rho,\eps}(K) e_n = q^{n+\eps} e_n, \qquad
\pi^{\mathrm P}_{\rho,\eps} \big(K^{-1}\big) e_n = q^{-(n+\eps)}e_n,\nonumber\\
\big(q^{-1}-q\big)\pi^{\mathrm P}_{\rho,\eps}(E) e_n = q^{-\frac12-n-\eps} \sqrt{(1-q^{2n+2\eps+2i\rho+1})
(1-q^{2n+2\eps-2i\rho+1})} e_{n+1},
\nonumber\\
\big(q^{-1}-q\big)\pi^{\mathrm P}_{\rho,\eps}(F) e_n = -q^{\frac12-n-\eps}
\sqrt{(1-q^{2n+2\eps+2i\rho-1}) (1-q^{2n+2\eps-2i\rho-1})} e_{n-1},
\nonumber\\
\big(q^{-1}-q\big)^2 \pi^{\mathrm P}_{\rho,\eps} (\Omega) e_n
=\big(q^{2i\rho}+q^{-2i\rho} -2\big) e_n.\label{eq:princ}
\end{gather}
For $(\rho,\eps)= (0,\frac12)$ the representation $\pi^{\mathrm P}_{0,\frac12}$ splits into a direct
sum of a positive and a negative discrete series representation: $\pi^{\mathrm P}_{0,\frac12}= \pi^+_\frac12
\oplus \pi^-_\frac12$. The representation space splits into two invariant subspaces: $\{ e_n\, | \, n \geq 0 \} \oplus \{e_n \, |\,
n<0 \}$.

\textbf{Complementary series.} The complementary series
representations $\pi^{\mathrm C}_{\lambda,\eps}$ are labeled by $\lambda$ and
$\eps$, where $\eps \in [0,\frac12)$ and $\lambda \in (-\frac12,-\eps)$, or
$\eps \in (\frac12,1)$ and $\lambda \in (-\frac12, \eps-1)$. The
representation space is $\ell^2(\Z)$ with orthonormal basis
$\{e_n\}_{n \in \Z}$. The action of the generators is the same as for the principal unitary series, with $i\rho$ replaced by $\lambda+\frac12$:
\begin{gather}
\pi^{\mathrm C}_{\lambda,\eps}(K) e_n = q^{n+\eps}\, e_n, \qquad
\pi^{\mathrm C}_{\lambda,\eps} \big(K^{-1}\big) e_n = q^{-(n+\eps)}e_n,\nonumber\\
\big(q^{-1}-q\big)\pi^{\mathrm C}_{\rho,\eps}(E) e_n
= q^{-\frac12-n-\eps} \sqrt{(1-q^{2n+2\eps+2\lambda+2})
(1-q^{2n+2\eps-2\lambda})} e_{n+1},
\nonumber\\
\big(q^{-1}-q\big)\pi^{\mathrm C}_{\lambda,\eps}(F) e_n = -q^{\frac12-n-\eps}
\sqrt{(1-q^{2n+2\eps+2\lambda}) (1-q^{2n+2\eps-2\lambda-2})} e_{n-1},\nonumber\\
\big(q^{-1}-q\big)^2 \pi^{\mathrm C}_{\lambda,\eps} (\Omega) e_n
=\big(q^{2\lambda+1}+q^{-1-2\lambda} -2\big) e_n.\label{eq:comp}
\end{gather}


\textbf{Strange series.} The f\/ifth class consists of the strange series representations $\pi^{\mathrm S}_{\sigma,\eps}$, labeled by $\sigma>0$ and $\eps \in [0,1)$, with representation space $\ell^2(\Z)$. The actions of the generators are the same as in principal unitary series with $i\rho$ replaced by $- \frac{ i\pi}{2\ln q}+\sigma$:
\begin{gather}
\pi^{\mathrm S}_{\sigma,\eps}(K) e_n = q^{n+\eps} e_n, \qquad
\pi^{\mathrm S}_{\sigma,\eps} \big(K^{-1}\big) e_n = q^{-(n+\eps)} e_n,\nonumber\\
\big(q^{-1}-q\big)\pi^{\mathrm S}_{\sigma,\eps}(E) e_n
= q^{-\frac12-n-\eps} \sqrt{(1+q^{2n+2\eps+2\sigma+1})
(1+q^{2n+2\eps-2\sigma+1})} e_{n+1},\nonumber\\
\big(q^{-1}-q\big)\pi^{\mathrm S}_{\sigma,\eps}(F) e_n = -q^{\frac12-n-\eps}
\sqrt{(1+q^{2n+2\eps+2\sigma-1}) (1+q^{2n+2\eps-2\sigma-1})} e_{n-1},\nonumber\\
\big(q^{-1}-q\big)^2 \pi^{\mathrm S}_{\sigma,\eps} (\Omega) e_n =-\big(q^{2\sigma}+q^{-2\sigma} +2\big)e_n.\label{eq:strange}
\end{gather}

The representations $\pi^{\mathrm P}_{\rho,\eps}$, $\pi^{\mathrm C}_{\lambda,\eps}$, $\pi^{\mathrm S}_{\sigma,\eps}$ with $\eps \in \R\setminus [0,1)$ are equivalent to the representations $\pi^{\mathrm P}_{\rho,\eps+m}$,  $\pi^{\mathrm C}_{\lambda,\eps+m}$, $\pi^{\mathrm S}_{\sigma,\eps+m}$ respectively, where $m \in \Z$ is such that $\eps+m \in [0,1)$. Therefore we will use these representations with any $\eps \in \R$. Furthermore, the principal unitary series representation~$\pi_{\rho,\eps}^{\mathrm P}$ is equivalent to $\pi_{-\rho,\eps}^{\mathrm P}$ and to $\pi_{\rho',\eps}^{\mathrm P}$ with $\rho'=\rho+\frac{m\pi}{\ln q}$, $m \in \Z$, and the complementary series representation $\pi_{\lambda,\eps}^{\mathrm{C}}$ is equivalent to $\pi_{-\lambda-1,\eps}^{\mathrm{C}}$.

It will be convenient to use the convention $e_{-n}=0$ for $n=1,2,\ldots$, if $e_n$ is a basis vector of $\ell^2(\N)$. Furthermore, we sometimes denote a basis vector of the representation space of the positive discrete series by $e_n^+$, and similarly for basis vectors of other representation spaces.


\section{A quantum analog of the tensor product of negative\\ and positive discrete series representations} \label{sec:-+}

Let $k_1,k_2>0$ and $\eps \in \mathbb R$. We consider the $\U_q$-representation
\begin{gather} \label{eq:defT1}
\mathcal T = \mathcal T_{k_1,k_2,\eps}= \Big(\big( \pi_{k_1}^- \otimes \pi_{k_2}^+\big) \oplus \big(\pi^{\mathrm S}_{k_1-\frac12,-\eps-k_1}  \otimes \pi_{k_2-\frac12,\eps+k_2}^{\mathrm S}\big)\Big)  \Delta.
\end{gather}
For the decomposition of $\mathcal T$ into irreducible representations we need the big $q$-Jacobi functions~\cite{KS03}. Let $a$, $b$, $c$ be parameters satisfying the conditions $a,b,c>0$ and $ab,ac,bc>1$. The big $q$-Jacobi functions are def\/ined by
\[
\Phi_\gamma(y;a,b,c|q)= \frac{ (\gamma/a, 1/bc, -y/abc\gamma;q)_\infty} { (1/ab, 1/ac, -y/bc;q)_\infty} \rphis{3}{2}{ 1/b\gamma, 1/c\gamma, -y/bc}{1/bc, -y/abc\gamma}{q, \gamma/a},
\]
for $y \in \C\setminus(-\infty, -bc]$ and $|\gamma|< a$. By \cite[(III.9)]{GRa04} $\Phi_\gamma$ is symmetric in $\gamma$, $\gamma^{-1}$, so for $|\gamma|\geq a$ we def\/ine the big $q$-Jacobi function by the same formula with $\gamma$ replaced by $\gamma^{-1}$. If $|y|<bc$, the $_3\varphi_2$-function can be transformed by \cite[(III.10)]{GRa04} and then we have
\begin{gather} \label{eq:3phi2}
\Phi_\gamma(y;a,b,c|q)= \rphis{3}{2}{\gamma/a, 1/a\gamma, -1/y}{1/ab, 1/ac}{q,-y/bc}.
\end{gather}
Let $t>0$. We def\/ine discrete sets $\Gamma^{\mathrm{f\/in}}$ and $\Gamma^{\mathrm{inf}}$ (`f\/in' and `inf' stand for `f\/inite' and `inf\/inite') by
\begin{gather*}
\Gamma^{\mathrm{f\/in}}=  \left\{ \frac{q^k}{e}\, \Big| \, e \in \{a,b,c\},\, k \in \N \text{ such that } \frac{q^k}{e}>1 \right\},\\
\Gamma^{\mathrm{inf}} = \left\{ -\frac{ abcq^k }{t}\, \Big| \, k \in \Z \text{ such that } \frac{abcq^k}{t}>1 \right\},
\end{gather*}
Note that the set $\Gamma^{\mathrm{f\/in}}$ is empty if $a,b,c>1$.
Now we introduce the measure $\nu(\,\cdot\,)=\nu(\,\cdot\,;a,b,c;t|q)$ by
\[
\int f(\gamma) d\nu(\gamma) = \frac1{4\pi i}\int_\T f(\gamma) v(\gamma) \frac{ d\gamma}{\gamma} + \sum_{\gamma \in \Gamma^{\mathrm{f\/in}} \cup \Gamma^{\mathrm{inf}}} f(\gamma)w(\gamma) ,
\]
where $v$ is the weight function on the counter-clockwise oriented unit circle $\T$ given by
\begin{gather*}
v(\gamma)=v(\gamma;a,b,c;t|q) =  \frac{ \theta(-t;q) }{\theta(-t/ab, -t/ac, -t/bc;q) (q,1/ab,1/ac,1/bc;q)_\infty}\nonumber\\
\phantom{v(\gamma)=v(\gamma;a,b,c;t|q) =}{}  \times \frac{ (\gamma^{\pm 2};q)_\infty }{ (\gamma^{\pm 1}/a, \gamma^{\pm 1}/b, \gamma^{\pm 1}/c;q)_\infty\, \theta(-t\gamma^{\pm 1}/abc;q)},%\label{eq:weight}
\end{gather*}
and{\samepage
\[
w(\gamma)=w(\gamma;a,b,c;t;q) = \Res{\gamma'= \gamma} \frac{ v(\gamma') }{\gamma'}, \qquad \gamma \in \Gamma^{\mathrm{f\/in}} \cup \Gamma^{\mathrm{inf}}.
\]
The weights $w$ in the discrete mass points can be calculated explicitly, see \cite[Section~8]{KS03}.}

Let $\mathcal H=\mathcal H(a,b,c;t|q)$ be the Hilbert space consisting of functions that satisfy $f(\gamma)=f(\gamma^{-1})$ ($\nu$-a.e.) with inner product
\[
\langle f ,g \rangle = \int f(\gamma) \overline{g(\gamma)}\, d\nu(\gamma).
\]
The set of functions $\{\gamma \mapsto \Phi_\gamma(x;a,b,c)\, | \, x \in \{-q^k\mid k \in \N\}\cup\{tq^k\mid k \in \Z \}\}$ is an orthogonal basis for $\mathcal H$, i.e.
\[
\langle \Phi_\cdot(y), \Phi_\cdot(y') \rangle = N_y \delta_{yy'}, \qquad y,y' \in \big\{{-}q^k\mid k \in \N\big\}\cup\big\{tq^k \mid k \in \Z\big\},
\]
where the quadratic norm $N_y$ is given by
\[
N_y =N_y(a,b,c;q)= \begin{cases}
 \dfrac{ (q,1/bc;q)_k }{(1/ab,1/ac;q)_k}q^{-k}, & y=-q^k, k \in \N ,\vspace{2mm}\\
 \dfrac{ (q,1/bc,-tq^k/ab, -tq^k/ac;q)_\infty }{ (1/ab,1/ac,-tq^k/bc, -tq^{k+1};q)_\infty }\frac{q^{-k}}{t}, & y=tq^k, k \in \Z.
\end{cases}
\]
We def\/ine
\[
\phi_y(x)=\phi_y(x;a,b,c|q) = \frac{\Phi_x(y;a,b,c|q)}{\sqrt{N_y(a,b,c;q)}},
\]
then $\{ \phi_{-q^k} \}_{k \in \N} \cup \{\phi_{tq^k} \}_{k \in \Z}$ is an orthonormal basis for $\mathcal H$. For ease of notations, it is useful to def\/ine $\phi_{-q^k}=0$ for $k=-1,-2,\ldots$.

The dif\/ference equation for the big $q$-Jacobi functions leads to a recurrence relation for the functions $\phi_{zq^k}$, $z\in\{-1,t\}$, of the form
\begin{gather} \label{eq:3termrecphi}
\big(x+x^{-1}\big)\phi_{zq^k}(x) = a_k \phi_{zq^{k+1}}(x) + b_k \phi_{zq^k}(x) + a_{k-1}\phi_{zq^{k-1}}(x),
\end{gather}
where
\begin{gather*}
a_k =a_k(a,b,c;z|q)=  \frac{ abcq^{-\frac12-2k} }{z^2} \sqrt{ \big(1+ zq^{k+1} \big) \left(1+\frac{ zq^k}{ab} \right) \left(1+\frac{zq^k}{ac} \right) \left(1+\frac{zq^k}{bc} \right) },\\
b_k =b_k(a,b,c;z|q)=  a+a^{-1} -\frac{abcq^{-2k}}{z^2}\left( 1+ \frac{zq^k}{ab} \right) \left( 1+ \frac{zq^k}{ac} \right)\\
\phantom{b_k =b_k(a,b,c;z|q)=}{}  -\frac{abcq^{1-2k}}{z^2} \big( 1+zq^k \big) \left( 1+\frac{zq^{k-1}}{bc} \right).
\end{gather*}
Here we assume $k \in \N$ if $z=-1$ and $k \in \Z$ if $z=t$, and we use $a_{-1}(a,b,c;-1|q)=0$. Note that $a_k$ and $b_k$ are both symmetric in $a,b,c$.

Let us remark that for $z=-1$ the recurrence relation corresponds to the three term recurrence relation for the continuous dual $q^{-1}$-Hahn polynomials, i.e.~Askey--Wilson polynomials with one parameter equal to zero and base $q^{-1}$;
\[
p_k(x)=p_k\big(x;a,b,c|q^{-1}\big) = a^{-k} \sqrt{ \frac{(ab,ac;q^{-1})_k}{(q^{-1}, bc;q^{-1})_k} } \rphis{3}{2}{q^{k}, ax, a/x}{ab, ac}{q^{-1}, q^{-1}}.
\]
So the function $\phi_{-q^k}$ is a multiple of the continuous dual $q^{-1}$-Hahn polynomials $p_k$. The moment problem corresponding to the continuous dual $q^{-1}$-Hahn polynomials is indeterminate, so the measure $\nu$ def\/ined above is a (non-extremal) solution for this moment problem.

We now turn to the decomposition of $\mathcal T$. First we diagonalize $\mathcal T(\Omega)$ on a suitable subspace of $\ell^2(\N)^{\otimes 2} \oplus \ell^2(\Z)^{\otimes 2}$. We def\/ine subspaces $\mathcal A_p$, $p\in\Z$, by
\[
\mathcal A_p =  \mathrm{span}\{f_{p,n}^-\ | \ n\in \N\}  \oplus  \mathrm{span} \{f_{p,n}^+\ | \ n \in \Z \},
\]
where
\[
f_{p,n}^- = \begin{cases}
e_n \otimes e_{n-p}, & p \leq 0,\\
e_{n+p} \otimes e_n, & p \geq 0,
\end{cases}
 \qquad
f_{p,n}^+=
\begin{cases}
(-1)^n e_{-n} \otimes e_{p+n}, & p \leq 0,\\
(-1)^{n+p} e_{-n-p} \otimes e_n, & p \geq 0.
\end{cases}
\]
It is useful to observe that $f_{p,n}^-=0$ for $n=-1,-2,\ldots$. Furthermore, note that $\overline{\mathcal A_p} \cong \ell^2(\N) \oplus \ell^2(\Z)$ and $\mathcal A = \bigoplus_{p \in \Z} \mathcal A_p$ is a dense subspace of $\ell^2(\N)^{\otimes 2} \oplus \ell^2(\Z)^{\otimes 2}$. From \eqref{eq:De(OM)} it follows that each subspace $\mathcal A_p$ is invariant under $\mathcal T(\Omega)$.

\begin{proposition} \label{prop:Te}
Let $p \in \Z$, $\sigma \in \{-,+\}$, and let $a$, $b$, $c$, $t_\sigma$ be defined by
\begin{gather}
a= \begin{cases}
q^{2k_2-2k_1-2p-1}, & p \geq 0,\\
q^{2k_1-2k_2+2p-1}, & p \leq 0,
\end{cases}
 \qquad b = q^{1-2k_1-2k_2}, \qquad
c=\begin{cases}
q^{2k_1-2k_2-1}, & p \geq 0,\\
q^{2k_2-2k_1-1}, & p \leq 0,
\end{cases}
\nonumber\\  t_\sigma =
\begin{cases}
-1, & \sigma=-,\\
q^{2\eps}, & \sigma=+.
\end{cases}\label{eq:abct}
\end{gather}
Then the operator $\Theta:\mathcal A_p \rightarrow \mathcal H(a,b,c;t_+|q^2)$ defined by
\[
f_{p,n}^\sigma \mapsto \left(x\mapsto \phi_{t_\sigma q^{2n}}\big(x;a,b,c|q^2\big) \right), \qquad n \in \Z,
\]
extends to a unitary operator. Moreover, $\Theta$ intertwines $\mathcal T(\Omega)$ with the multiplication operator $(q^{-1}-q)^{-2}M_{x+x^{-1}-2}$.
\end{proposition}

\begin{proof}
First of all, $\Theta$ extends to a unitary operator since it maps one orthonormal basis onto another.

To check the intertwining property, we consider the explicit action of $\Omega$.
First assume $p \geq 0$. From \eqref{eq:De(OM)}, \eqref{eq:pos}, \eqref{eq:neg} and \eqref{eq:strange} it follows that the action of the Casimir operator is given by
\[
\big(q^{-1}-q\big)^2 [\mathcal T(\Omega)+2 ] f_{p,n}^\sigma = a_n^\sigma f_{p,n+1}^\sigma + b_n^\sigma f_{p,n}^\sigma + a_{n-1}^\sigma f_{p,n-1}^\sigma,
\]
where
\begin{gather*}
a_n^\sigma =   t_\sigma^{-2}q^{-2-2k_2-2k_1-2p-4n}\\
 \hphantom{a_n^\sigma =}{}
 \times \sqrt{ \big(1+t_\sigma q^{4k_1+2p+2n}\big)\big(1+t_\sigma q^{2p+2n+2}\big)\big(1+t_\sigma q^{4k_2+2n}\big) \big(1+t_\sigma q^{2n+2}\big)},\\
b_n^\sigma=   q^{-2p-2k_1+2k_2-1} +q^{2p+2k_1-2k_2+1} - t_\sigma^{-2}q^{1-2k_1-2k_2-2p-4n}\big(1+t_\sigma q^{2n+4k_2-2}\big)\big(1+t_\sigma q^{2n}\big) \\
\phantom{b_n^\sigma=}{}
 + t_\sigma^{-2}q^{-1-2k_1-2k_2-2p-4n}\big(1+t_\sigma q^{4k_1+2p+2n}\big)\big(1+t_\sigma q^{2p+2n+2}\big).
\end{gather*}
The intertwining property follows from comparing this with the recurrence relation \eqref{eq:3termrecphi} for the function $\phi_{t_\sigma q^{2n}}(x;a,b,c|q^2)$. For $p <0$ the proof runs along the same lines.
\end{proof}

By Proposition~\ref{prop:Te} $\Theta$ intertwines $\mathcal T(\Omega)$ with a multiplication operator on~$\mathcal H$. This allows us to read of\/f the spectrum of $\mathcal T(\Omega)$ from the support of the measure $\nu$. Since the irreducible $*$-representations are characterized by the actions of $\Omega$ and $K$, we now only need to consider the action of $K$ to f\/ind the decomposition of $\mathcal T$ into irreducible representations;
\[
\mathcal T(K) f_{p,n}^\sigma = q^{-2p + 2k_2 -2k_1}  f_{p,n}^\sigma.
\]
By comparing the spectral values of $\mathcal T(K)$ and $\mathcal T(\Omega)$ with \eqref{eq:pos}--\eqref{eq:strange}, we obtain the following decomposition of $\mathcal T$.

\begin{theorem} \label{thm:decompT1}
The $\U_q$-representation $\mathcal T_{k_1,k_2,\eps}$ is unitarily equivalent to:
\begin{alignat*}{4}
& (i)\quad && \dirint{0}{-\frac{\pi}{2 \ln q}} \pi^{\mathrm P}_{\rho,\eps'} d\rho \oplus \bigoplus_{\substack{j \in \Z \\ \sigma_j>0}} \pi^{\mathrm S}_{\sigma_j,\eps'},\qquad  &&  |k_1-k_2|\leq \frac12,\ k_1+k_2\geq \frac12,& \\
% Original source lines 382--385
\end{alignat*}
where $\eps'=k_2-k_1$, $\sigma_j= k_1+k_2+\eps+\frac12+j$, $k_j^+=k_2-k_1-j$ and $k_j^-=k_1-k_2-j$. For $n,p \in \Z$ and $\sigma \in \{-,+\}$ the unitary intertwiner is given by
\begin{alignat*}{3}
& (i)\quad && f_{p,n}^\sigma \mapsto \Lambda f_{p,n}^\sigma,  & \\
% Original source lines 398--405
\end{alignat*}
where
\begin{gather*}
\Lambda f_{p,n}^\sigma =  \int_0^{-\frac{\pi}{2\ln q}} \big(\Theta f_{p,n}^\sigma\big)\big(q^{2i\rho}\big) \sqrt{v\big(q^{2i\rho};a,b,c;t_+|q^2\big)}  e_{-p}^{\mathrm P}\, d\rho\\
\phantom{\Lambda f_{p,n}^\sigma =}{} + \sum_{\substack{j \in \Z\\\sigma_j>0}} \big(\Theta f_{p,n}^\sigma\big)\big({-}q^{2\sigma_j}\big) \sqrt{w\big({-}q^{2\sigma_j};a,b,c;t_+;q^2\big)} e_{-p}^{\mathrm S},
\end{gather*}
$\Theta$ is as in Proposition {\rm \ref{prop:Te}} and $a$, $b$, $c$, $t_+$ are given by \eqref{eq:abct}.
\end{theorem}
\setcounter{section}{3}
\setcounter{subsection}{0}
\setcounter{subsubsection}{0}
\setcounter{equation}{4}
% Original source lines 420--430
Let us denote the intertwiner from Theorem \ref{thm:decompT1} by $I$. The actions $I\circ \mathcal T(X)\circ I^{-1}$, $X=E,F$, are given by the appropriate actions of $E$ and $F$ in \eqref{eq:pos}--\eqref{eq:strange}, up to a phase factor. It is possible to determine the actions explicitly using the explicit expressions for the weight functions~$v$ and~$w$, the explicit expressions for $\Theta f_{p,n}^\sigma$ as $_3\varphi_2$-functions, and the following contiguous relations for $_3\varphi_2$-functions, see \cite[(2.3), (2.4)]{GIM96},
\begin{gather*}
\rphis{3}{2}{Aq,B,C}{D,E}{q,\frac{DE}{ABCq}} - \rphis{3}{2}{A,B,C}{D,E}{q,\frac{DE}{ABC}}  \\
\qquad{}=\frac{DE(1-B)(1-C)}{ABCq(1-D)(1-E)} \rphis{3}{2}{Aq,Bq,Cq}{Dq,Eq}{q,\frac{DE}{ABCq}},
\\
\left(1-\frac{D}{A}\right)\left(1-\frac{E}{A}\right) \rphis{3}{2}{A/q,B,C}{D,E}{q,\frac{DEq}{ABC}} \\
\qquad{} -\left(1-\frac{q}{A}\right)\left(1-\frac{DE}{ABC}\right)\rphis{3}{2}{A,B,C}{D,E}{q,\frac{DE}{ABC}}
\\
\qquad{}=\frac{q}{A}\left(1-\frac{D}{q} \right) \left(1-\frac{E}{q} \right)\rphis{3}{2}{A/q,B/q,C/q}{D/q,E/q}{q,\frac{DEq}{ABC}}.
\end{gather*}
We do not work out the details.
\begin{thebibliography}{99}
\bibitem[2]{BK93}
Burban I.M., Klimyk A.U.,
Representations of the quantum algebra $U_q(su_{1,1})$,
\href{http://dx.doi.org/10.1088/0305-4470/26/9/011}{{\it J.~Phys.~A: Math.~Gen.}} \textbf{26} (1993), 2139--2151.

\bibitem[3]{GRa04}
Gasper G., Rahman M.,
Basic hypergeometric series, 2nd ed.,
\href{http://dx.doi.org/10.1017/CBO9780511526251}{{\it Encyclopedia of Mathematics and its Applications}}, Vol.~96, Cambridge University Press, Cambridge, 2004.

\bibitem[9]{GIM96}
Gupta D.P., Ismail M.E.H., Masson D.R.,
Contiguous relations, basic hypergeometric functions, and ortho\-go\-nal polynomials. III.~Associated continuous dual $q$-Hahn polynomials,
\href{http://dx.doi.org/10.1016/0377-0427(95)00264-2}{{\it J.~Comput. Appl. Math.}} \textbf{68} (1996), 115--149,
\href{http://arxiv.org/abs/math.CA/9411226}{math.CA/9411226}.

\bibitem[13]{KS03}
Koelink E., Stokman J.V.,
The big $q$-Jacobi function transform,
\href{http://dx.doi.org/10.1007/s00365-002-0498-x}{{\it Constr. Approx.}} \textbf{19} (2003), 191--235,
\href{http://arxiv.org/abs/math.CA/9904111}{math.CA/9904111}.

\bibitem[15]{MMNNSU}
Masuda T., Mimachi K., Nakagami Y., Noumi M., Saburi Y., Ueno K.,
Unitary representations of the quantum group ${\rm SU}_q(1,1)$. II.~Matrix elements of unitary representations and the basic hypergeometric functions,
\href{http://dx.doi.org/10.1007/BF01039312}{{\it Lett. Math. Phys.}} {\bf 19} (1990),  195--204.

\bibitem[20]{Sch}
Schm\"udgen K.,
Unbounded operator algebras and representation theory,
{\it Operator Theory: Advances and Applications}, Vol.~37, Birkh\"auser Verlag, Basel, 1990.

\end{thebibliography}
\end{document}
